% model_specification.tex -- Slide-to-Target: model specification (reference note, v1, 31 Oct 2026).
%
% A reference note for the team, written alongside the code (physics.py, observe.py, belief.py, decide.py,
% drake_config.py, e0_*.py).  It replaces the long-form proposal's paragraphs "Assumptions on the object",
% "Contact model" and "Fusing the prior with practice slides".  Under the 6.4210 AI policy the team writes the
% prose of the submitted proposal and report; this note supplies the equations, tables and figures around it.
% Build: pdflatex model_specification.tex (twice).  error_budget_rows.tex is generated by `python3 observe.py`.
\documentclass[10pt]{article}
\usepackage[utf8]{inputenc}
\usepackage[T1]{fontenc}
\usepackage[margin=2.2cm]{geometry}
\usepackage{amsmath,amssymb}
\usepackage{array,booktabs,longtable}
\usepackage{hyperref}
\usepackage{xcolor}
\setlength{\parskip}{4pt}
\setlength{\parindent}{0pt}
\newcommand{\dstar}{d^{\star}}
\newcommand{\code}[1]{\texttt{#1}}

\title{Slide-to-Target: model specification\\[2pt]\large reference note v1}
\author{Slide-to-Target team, MIT 6.4210}
\date{31 October 2026}

\begin{document}
\maketitle

\noindent\fbox{\parbox{0.97\textwidth}{\small
\textbf{What this is.} The full generative model the instructor asked for: how a commanded push becomes a release
state, how a release state becomes a landing, what is measured, and how the belief over friction is updated and
used.  Sections~\ref{sec:state}--\ref{sec:ladder} are the specification; Section~\ref{sec:impl} maps every equation
to a function and a test in the repository.  Numbers marked \emph{provisional} wait on the E0 runs named beside them.
\textbf{Not submitted text.} The team writes the proposal and report prose; this note is the reference the code is
checked against.}}

\section{Overview}
\label{sec:overview}
The pipeline, one practice slide per loop:
\begin{center}\small
\begin{tabular}{@{}l@{\;$\rightarrow$\;}l@{\;$\rightarrow$\;}l@{\;$\rightarrow$\;}l@{}}
commanded push $u$ & push map (release model) & free slide & measurements: frames, rest pose, FK release point\\[2pt]
\multicolumn{4}{@{}l}{$\downarrow$ physics likelihood (model C) \quad $\leftarrow$ LLM observation model, prior}\\[2pt]
\multicolumn{4}{@{}l}{grid posterior over $\theta=\ln\mu$ \;$\rightarrow$\; decision (minimum expected loss) \;$\rightarrow$\; next push}
\end{tabular}
\end{center}
The first shot uses the LLM and the prior alone; each practice slide closes the loop once.  Everything below
holds for the main study; the fidelity ladder of Section~\ref{sec:ladder} says which assumptions are relaxed
where.

\section{State and parameters}
\label{sec:state}
\begin{itemize}
\item \textbf{Unknown, shared across a pair's slides:} $\theta=\ln\bar\mu$, the mean object--table kinetic friction,
      and strip offsets $\eta_s$ (level L2 and up).
\item \textbf{Known and fixed in the main study:} mass $m$, geometry, centre of mass at the base centroid (offset $e=0$),
      paddle--object friction $\mu_p$.
\item \textbf{Free at level L3:} $e$ and $\mu_p$.
\item \textbf{Per-slide nuisance:} release position $x_r$ and release speed $v_r$.  They are estimated from data, never
      taken from the commanded trajectory.
\end{itemize}
Under A1--A3 the stopping distance does not depend on the centre of mass: pure translation puts the friction
resultant under the COM, and a line contact can place its force centroid anywhere along the contact segment.  So
the COM is treated as known for boxes, and E0-3 tests where that fails.

\section{Assumptions, each with its check}
\label{sec:assumptions}
\begin{center}\small
\begin{tabular}{@{}llp{3.2cm}@{}}
\toprule
ID & Assumption & Checked by\\
\midrule
A1 & Flat rigid base, uniform pressure & E0-1\\
A2 & Line contact; the COM projection lies inside the paddle contact segment ($|e|<w/2$) & E0-3\\
A3 & Push along the face normal; separation at the first time the paddle decelerates faster than $\mu g$ & E0-4\\
A4 & Friction strips run across the push direction; $\mu$ does not depend on speed over 0.3--1.2 m/s & E0-6; rig $\mu(v)$\\
A5 & No tipping during the push, $(a+\mu g)h_p\le gb+a\,h_c$, or during the slide, $\mu h_c<b$ & analytic, per object\\
A6 & Drake: 1 ms step, Lagged approximation (tentative), one compliant geometry per contact pair, margin $10^{-4}$ m & E0-1, E0-2\\
\bottomrule
\end{tabular}
\end{center}
In A5, $b$ is the half-length along the push, $h_c$ the COM height and $h_p$ the paddle contact height
(\code{physics.no\_tip}).

\section{Contact specification}
\label{sec:contact}
\begin{center}\small
\begin{tabular}{@{}p{3.1cm}p{6.2cm}p{6.2cm}@{}}
\toprule
Property & Paddle--object & Object--table\\
\midrule
Geometry & Flat paddle face on the box's end face & Box base on the table plane\\
Contact type in the plane & Segment of width $w$ (line).  Discs in E0-3 give point contact on purpose & Full base patch\\
Dynamic friction & $\mu_p$ chosen (nominal 0.3), swept 0.1--0.6 in E0-3 & Measured on the rig (primary set), or Yu et al.\ (simulator-only set)\\
Static friction & Not used by Drake's discrete solver & Not used\\
Compliance & Object compliant, paddle rigid & Object compliant, table rigid\\
Hydroelastic modulus & $10^{7}$ Pa nominal (Drake default); swept $10^{5}$--$10^{7}$ & Same\\
Dissipation & Hunt--Crossley $d$ under Lagged/Similar (swept 0, 1, 10, 50 s/m); \code{relaxation\_time} only under SAP & Same\\
Margin & $10^{-4}$ m & $10^{-4}$ m\\
Contact height & $h_p$ chosen to satisfy A5 & Not applicable\\
How it affects the landing & Through release speed and yaw (E0-3, E0-4) & Directly (E0-1)\\
\bottomrule
\end{tabular}
\end{center}
Two Drake facts shape this table.  With two rigid hydroelastic geometries, the default
\code{hydroelastic\_with\_fallback} silently switches that pair to point contact; and if mesh elements are much
larger than the contact surface, contact effectively becomes point contact.  Experiments therefore run the strict
\code{hydroelastic} model, which throws on a rigid--rigid pair (\code{drake\_config.py}, E0-2).  At
$E=10^{7}$ Pa the paddle contact penetrates about 2.5 $\mu$m, so a $0.002^\circ$ yaw misalignment puts the whole
patch at one edge: the paddle contact is modelled as a segment whose force centroid can move (Lynch and Mason), not
as an averaged face.  Drake combines a pair's friction as $2\mu_a\mu_b/(\mu_a+\mu_b)$; the object geometry carries a
fixed value and each strip and the paddle get the partner value that makes the pair coefficient equal to its target
(\code{scene.partner\_mu}), so one object can touch strips of different friction.

\section{Generative model}
\label{sec:generative}
\paragraph{Push map (release model).}  Contact holds while the paddle decelerates more slowly than friction would
decelerate the free object:
\begin{equation}
t^{*}=\min\{\,t>t_{\text{peak}}:\ a_p(t)<-\mu g\,\},\qquad v_r=v_p(t^{*}).
\label{eq:release}
\end{equation}
For braking that starts at zero deceleration and ramps at constant jerk $J$,
\begin{equation}
v_r=v_{\text{peak}}-\frac{(\mu g)^2}{2J}.
\label{eq:vr}
\end{equation}
At $J=100$ m/s$^3$ and $v_{\text{peak}}=0.8$ m/s this costs 11\% of the distance at $\mu=0.3$ and about 1\% at
$\mu=0.1$ (\code{physics.release\_speed}; \code{release\_from\_profile} evaluates~\eqref{eq:release} on any sampled
paddle profile).  Under A2--A3, heading $\psi_r\approx0$ and spin $\omega_r\approx0$ at release.

\paragraph{Free slide.}  With uniform friction,
\begin{equation}
x(\tau)=x_r+v_r\tau-\tfrac12\mu g\tau^{2},\quad 0\le\tau\le T=\frac{v_r}{\mu g},\qquad
x_{\text{rest}}=x_r+\frac{v_r^{2}}{2\mu g}.
\label{eq:slide}
\end{equation}
Across strips the deceleration is $\mu(x)g$, integrated piecewise until the speed reaches zero
(\code{physics.slide\_piecewise}, \code{slide\_field}).  Spin is a diagnostic only.  For a base with symmetric
pressure spin \emph{lengthens} the slide, because sliding and spinning share the Coulomb friction (Farkas et al.):
a uniform-pressure disk gives $d/d_0\approx1+0.26\,(R\omega_r/v_r)^2$ for $R\omega_r/v_r\le1$, with 1.065 at 0.5 and
1.27 at 1 from the integrated model (\code{physics.disk\_spin\_slide}).  A heading error $\psi$ shortens the
along-axis distance only by $\psi^2/2$ and adds a lateral miss of $\psi d$.

\paragraph{Transition across practice slides.}  The parameters do not change between slides:
\begin{equation}
\phi_{j+1}=\phi_j .
\end{equation}
Yu et al.\ measured plywood friction falling from 0.28 to 0.24 over 100 passes, about 0.15\% per pass, which is
negligible over five slides, and the object is reset between slides.  This is a decision, not an omission.

\section{Observation models}
\label{sec:observation}
\paragraph{$O_{\text{LLM}}$: the language model as a sensor.}  The prior $p_0(\theta)$ is the population
distribution of measured $\ln\mu$, fitted on training folds (\code{belief.prior\_from\_measured}).  The LLM's
median $m$ enters through a likelihood whose parameters are fitted leave-one-object-out on measured pairs
(\code{belief.fit\_llm\_likelihood\_loo}):
\begin{equation}
m_{\text{LLM}}=\alpha+\beta\theta+\varepsilon,\qquad\varepsilon\sim\mathcal N(0,s^{2}).
\end{equation}
This is the old ``recalibration'' step, now the observation model.  If $\beta\to0$ the LLM carries no
information.  The prompt builder reads only the \code{description\_*} and \code{photo\_*} columns of the pair file
(\code{llm\_prior.py}, test 7); prompts are frozen and responses cached before the real measured file exists.

\paragraph{$O_{\text{phys}}$, model C (primary).}  Each practice slide gives three measurements of~\eqref{eq:slide}:
frames $\hat x(t_i)$ of the object during the free slide from the depth camera and ICP, noise $\sigma_f$ per frame;
the rest pose, averaged over many stationary frames, noise $\sigma_{\text{rest}}$; and the release point, from arm
forward kinematics at separation, noise $\sigma_{\text{FK}}$:
\begin{equation}
\hat x_i=x(t_i;x_r,v_r,\mu)+e_i,\qquad
\hat x_{\text{rest}}=x_r+\frac{v_r^{2}}{2\mu g}+e_{\text{rest}},\qquad
\hat x_r^{\text{FK}}=x_r+e_{\text{FK}} .
\end{equation}
Model A uses frames only; model B adds the rest pose; model C adds the release point
(\code{observe.fit\_slide}).  Frames are timed from the arm's separation time, $t_i=i/f$ for $0\le t_i\le T$.

\paragraph{O1 (ablation only).}  Endpoint plus commanded speed, $y=\ln\big(v_{\text{cmd}}^2/2gd\big)$.  It carries the
release-ramp bias $(v_{\text{cmd}}/v_r)^2-1$, which grows with $\mu$ by~\eqref{eq:vr}.  It is kept only to show that
bias (\code{observe.fit\_o1}, \code{o1\_bias}).

\paragraph{Error budget.}  Relative standard error of $\mu$ from one slide (Cram\'er--Rao bound,
\code{observe.crlb}; test 2 reproduces every cell).  Camera at 30 Hz with 2 mm per-frame noise; rest pose and FK
at 0.5 mm.  O1 assumes 3\% speed error and 2 mm endpoint error; its bias is computed at braking jerk $J$.
\begin{center}\small
\begin{tabular}{@{}cclrrrrrr@{}}
\toprule
$\mu$ & $v_r$ (m/s) & Slide & A: frames & B: $+$rest & C: $+$FK & O1 noise & O1 bias, $J{=}100$ & O1 bias, $J{=}300$\\
\midrule
% generated by observe.py: relative standard error of mu from one slide (30 Hz, 2 mm frames; 0.5 mm rest and FK)
0.1 & 0.8 & 0.82 s, 32.6 cm & 1.6\% & 1.1\% & 0.8\% & 6.0\% & +1.2\% & +0.4\% \\
0.1 & 1.1 & 1.12 s, 61.7 cm & 0.7\% & 0.5\% & 0.4\% & 6.0\% & +0.9\% & +0.3\% \\
0.3 & 0.5 & 0.17 s, 4.2 cm & 20.0\% & 15.6\% & 12.1\% & 7.6\% & +19.9\% & +6.0\% \\
0.3 & 0.8 & 0.27 s, 10.9 cm & 7.0\% & 5.2\% & 3.9\% & 6.3\% & +11.8\% & +3.7\% \\
0.3 & 1.1 & 0.37 s, 20.6 cm & 3.3\% & 2.5\% & 1.8\% & 6.1\% & +8.4\% & +2.7\% \\
0.5 & 0.5 & 0.10 s, 2.5 cm & 36.7\% & 30.6\% & 25.6\% & 9.9\% & +73.4\% & +18.2\% \\
0.5 & 0.8 & 0.16 s, 6.5 cm & 19.6\% & 10.3\% & 8.0\% & 6.7\% & +38.5\% & +10.8\% \\
0.5 & 1.1 & 0.22 s, 12.3 cm & 8.0\% & 4.9\% & 3.7\% & 6.2\% & +26.1\% & +7.7\% \\

\bottomrule
\end{tabular}
\end{center}
Two consequences.  Model C is unbiased but its variance grows fast as slides get short: the error scales roughly
as $\sigma_f/(\sqrt N\,T^{2})$ with $T=v/\mu g$.  And the prior's worth in practice slides,
$n_0=\sigma_{\text{obs}}^{2}/s_0^{2}$ (\code{belief.n0}), now depends on the object: with $s_0=0.45$ it ranges from
about $10^{-4}$ slides for slippery objects to about 0.3 for $\mu=0.5$ at 0.5 m/s.

\section{Inference}
\label{sec:inference}
A grid of 500 points over $\theta\in[\ln0.05,\ln1.0]$ (\code{belief.theta\_grid}).  The physics likelihood is
Gaussian in $\mu$ while the prior and the LLM model are Gaussian in $\ln\mu$, so the grid is exact where a conjugate
update is not (test 4 checks it against the conjugate closed form for Gaussian-in-$\theta$ observations).  The
per-slide nuisance $(x_r,v_r)$ is profiled out by least squares at each grid point
(\code{belief.profile\_nuisance}, vectorized Gauss--Newton):
\begin{equation}
p(\theta\mid m,z_{1:k})\;\propto\;p_0(\theta)\,p(m\mid\theta)\prod_{j=1}^{k}\max_{x_r,v_r}p(z_j\mid\theta,x_r,v_r).
\end{equation}
At level L3 the grid becomes two-dimensional over $(\mu,e)$.  Pushing outcomes can be heteroscedastic and
multimodal (Bauza), which a single Gaussian cannot represent but a grid can.  The grid's floor at $\mu=0.05$ acts as a
hard prior: for a wide posterior near it, the $\mu^{-2}$ weight of the squared-relative loss makes the decision
depend on that floor (test 5 notes the effect).

\section{Decision rule}
\label{sec:decision}
Because $v_r$ depends on $\mu$, the closed-form release speed only holds when braking is abrupt.  The commanded peak
speed $u$ is chosen by a one-dimensional numerical search on the expected loss, with the release model inside $d$
(\code{decide.choose\_speed}):
\begin{equation}
u^{*}=\arg\min_{u}\int L\big(d(\mu,u)-\dstar\big)\,p(\mu\mid\text{data})\,d\mu,\qquad
d(\mu,u)=\frac{v_r(\mu,u)^{2}}{2\mu g}.
\end{equation}
The old closed form, $v_{\text{opt}}^{2}=2g\,\dstar\exp(m_k-1.5S_k^{2})$, is kept as the unit-test limit for
$J\to\infty$ with squared relative loss (test 5).  Losses: squared relative error, absolute relative error, a miss
outside $\max(0.1\dstar,\,1\text{ cm})$, and the table-edge penalty at $L=1.3\dstar$.  With $\ln\mu\sim\mathcal N(\ln0.3,0.3^2)$
and $\dstar=0.3$ m, the optimum moves from 1.242 m/s at $J=\infty$ to 1.253 at $J=300$ and 1.274 at $J=100$ m/s$^3$.

\textbf{Practice slides run at the maximum feasible speed} (\code{decide.practice\_speed}).  The standard error
falls roughly as $v^{-2.5}$, so the Fisher information grows roughly as $v^{5}$.  That is valid only if $\mu$ does not
depend on speed, which the rig's $\mu(v)$ data check (A4).

\section{Fidelity ladder}
\label{sec:ladder}
\begin{center}\small
\begin{tabular}{@{}lp{3.0cm}p{2.9cm}p{2.6cm}p{2.2cm}p{4.0cm}@{}}
\toprule
Level & Unknowns & Observation & Inference & Experiments & Expected difficulty\\
\midrule
L0 & $\mu$ only, perfect measurements & $d$ and $v_r$ exact & closed form & pipeline check & trivial: one slide fixes $\mu$; the prior matters only at $k=0$\\
L1 & $+$ execution and perception noise & model C vs O1 & 1-D grid & E0-5, E2--E4 & O1's release-ramp bias; model C's variance at high $\mu$\\
L2 & $+$ friction field (strips) & model C per lane & grid with lane offsets & E0-6, E4 & which lane to practise on\\
L3 & $+$ COM offset $e$, point contact & $+$ landing yaw and lateral offset & 2-D grid $(\mu,e)$ & E0-3 & mostly lateral error; distance error second order\\
L4 & $+$ simulator vs analytic mismatch & as above & Drake rollouts as the likelihood & E0-1; Isaac branch & contact parameters, discrete approximation\\
\bottomrule
\end{tabular}
\end{center}
The residual of the constant-$\mu$ fit tells which level is active (\code{observe.residual\_signature}, by BIC over
closed-form alternatives): a clean parabola with residuals at noise level (constant $\mu$ on a uniform table); a
step in deceleration (a friction strip crossed); deceleration that changes smoothly (spin, as the spin-to-slide
ratio evolves; or a speed-dependent $\mu$, which across slides shows as $\hat\mu$ correlated with $v_r$,
\code{observe.speed\_signature}).  \emph{Detectability}: on synthetic 1.1 m/s slides with $\mu$ stepping from 0.1
to 0.2 part way, the step is labelled on 20 of 20 trials at 0.5 mm frame noise and 60 Hz and on 18 of 20 at 1 mm,
but at the budget's 2 mm and 30 Hz it reads as a smooth trend on 20 of 20.  Telling a strip from spin in one slide
needs sub-millimetre frames or 60 Hz and up, which bears on the camera-rate decision (E0-5).

Prediction to report against: L0 is easy, as the instructor expects.  The real difficulty is the $\mu$-dependent
release and the simulator-vs-analytic mismatch, not friction or COM uncertainty as such.

\section{Corrections to earlier drafts}
\label{sec:corrections}
\begin{enumerate}\small
\item Off-centre friction does not make the path ``curve and shorten'': for a symmetric base, spin lengthens the slide
      ($d/d_0\approx1+0.26(R\omega/v)^2$); a heading error $\psi$ shortens the along-axis distance only by $\psi^2/2$ and
      adds a lateral miss $\psi d$.
\item A trajectory fit estimates $\mu$ to 1.7\% only for long slides; at $\mu=0.5$, $v=0.5$ m/s frames alone give 37\%.
\item The standard error falls as $v^{-2.5}$, so the information grows as $v^{5}$ (not $v^{2.5}$).
\item SAP ignores Hunt--Crossley dissipation and uses \code{relaxation\_time}; Lagged and Similar use Hunt--Crossley
      dissipation and ignore \code{relaxation\_time}.  Each parameter is paired with its approximation.
\item Release speed is not the paddle's peak speed: $v_r=v_{\text{peak}}-(\mu g)^2/2J$, and the error grows with $\mu$.
\item Lookup rows for Yu's pairs are not Yu's values: Yu sets simulator truth only (\code{ground\_truth\_published.csv});
      the lookup (\code{friction\_lookup.csv} v2) carries independent values for the same pair IDs.
\item The rig tests objects O1--O9 on surfaces S1--S3 as defined in protocol v3, not the materials listed in lookup
      table v1.
\end{enumerate}

\section{Implementation map and status}
\label{sec:impl}
\begin{center}\small
\begin{tabular}{@{}lp{7.6cm}p{5.2cm}@{}}
\toprule
Module & Functions & Tests (\code{pytest tests/})\\
\midrule
\code{physics.py} & \code{stop\_distance}, \code{release\_speed}, \code{release\_from\_profile}, \code{slide\_trajectory}, \code{slide\_piecewise}, \code{slide\_field}, \code{spin\_factor}, \code{disk\_spin\_slide}, \code{no\_tip} & 1: sampled release vs~\eqref{eq:vr} to $10^{-6}$ m/s; 6: disk $d/d_0$ 1.065, 1.27\\
\code{observe.py} & \code{simulate\_obs}, \code{fit\_slide} (A, B, C), \code{fit\_o1}, \code{crlb}, \code{error\_budget}, \code{residual\_signature}, \code{speed\_signature} & 2: budget table to 0.1 pp; 3: 2,000 slides, bias $<1\%$, SD within 20\% of CRLB\\
\code{belief.py} & \code{theta\_grid}, \code{prior\_from\_measured}, \code{fit\_llm\_likelihood(\_loo)}, \code{loglik\_slide}, \code{posterior}, \code{n0} & 4: grid vs conjugate update to $10^{-3}$\\
\code{decide.py} & \code{expected\_loss}, \code{choose\_speed}, losses, \code{closed\_form\_speed}, \code{practice\_speed} & 5: $J\to\infty$ limit to three significant digits\\
\code{llm\_prior.py} & \code{build\_prompt}, \code{cached\_response} & 7: only \code{description\_*}, \code{photo\_*} reach a prompt\\
\code{drake\_config.py}, \code{sim\_config.yaml} & \code{load\_config}, \code{contact\_from\_config}, \code{describe}, \code{assert\_hydro\_only} & E0-2 (\code{e0\_2\_hydro\_check.py})\\
\code{e0\_1\_free\_slide.py} & grid, spin and heading launches & E0-1 pass criteria; spin within 10\% of the disk model\\
\bottomrule
\end{tabular}
\end{center}
\paragraph{E0 status (31 Oct).}  \emph{E0-2 passes}: in the passive and the push scene, for box, disc, ring and
footed objects, in strict and in fallback mode, every contact over 300 steps is hydroelastic (object--table every
step; paddle--object every step in the push scene) and there are zero point-pair contacts
(\code{e0\_2\_results.csv}, \code{e0\_2\_contacts.png}).  Negative controls: the paddle lowered 1 mm into the table
makes strict mode throw; a free rigid box on the rigid table runs silently on point contact in fallback mode.
\emph{E0-1 freezes A6} at the \code{sim\_config.yaml} values (Lagged, 1 ms, $E=10^{7}$ Pa, dissipation 1 s/m,
margin $10^{-4}$ m): box slides within 0.96\% of $d_0$ and $0.00^\circ$ of yaw over $\mu\in\{0.1,0.3,0.5\}$,
$v\in\{0.6,1.0\}$ m/s; disc spin launches within 2\% of the disk integrator (1.053--1.060 at $R\omega/v=0.5$,
1.249--1.266 at 1); heading launches within 1\% along and across with no yaw drift.  Finding: with the $10^{-4}$ m
margin, Hunt--Crossley dissipation of 10 s/m or more under Lagged (1 s/m or more under Similar) shortens slides by 4
to 85\%, while without margin Lagged passes up to 50 s/m; 2 ms steps pass only without margin, 5 ms never; SAP
passes at 1 ms with relaxation time 0.01 s or less (\code{e0\_1\_summary.txt}).  19 of 66 settings pass.

\appendix
\section{sim\_config.yaml}
\begin{verbatim}
# sim_config.yaml -- the one Drake contact configuration that every experiment loads (drake_config.py).
# Status: FROZEN by E0-1 on 31 Oct 2026 (e0_1_free_slide.py; e0_1_summary.txt, e0_1_summary.png).  Change only with a new E0-1 run.
#   E0-1 result at these settings: box free slide within 0.96% of d = v^2/(2 mu g) and 0.00 deg of yaw over
#   mu in {0.1, 0.3, 0.5}, v in {0.6, 1.0} m/s; disc spin launches within 2% of the disk model; heading launches pass.
#   Interaction to respect: with margin 1e-4 m, Hunt-Crossley dissipation >= 10 s/m (lagged) or >= 1 s/m (similar)
#   shortens slides by 4-85%; without margin, lagged passes up to 50 s/m.  Keep dissipation <= 1 s/m with this margin.
#   2 ms steps pass only without margin (lagged) and 5 ms steps never pass.  SAP passes at 1 ms with relaxation <= 0.01 s.
time_step: 0.001                         # s                          (A6)
discrete_contact_approximation: lagged   # lagged | similar | sap     (A6, tentative)
contact_model: hydroelastic              # strict: a rigid-rigid pair throws instead of silently falling back to point contact
hydroelastic_modulus: 1.0e7              # Pa, object (compliant); Drake's default; swept 1e5-1e7 in e0_sweep.py
dissipation: 1.0                         # s/m, Hunt-Crossley: used by lagged and similar, IGNORED by sap
relaxation_time: null                    # s, sap only (null = Drake's default); IGNORED by lagged and similar
margin: 1.0e-4                           # m, hydroelastic margin on every geometry
resolution_hint: 0.01                    # m, object mesh resolution
compliance:                              # one compliant geometry per contact pair
  object: compliant
  table: rigid
  paddle: rigid
paddle_mu: 0.3                           # paddle-object pair coefficient; swept 0.1-0.6 in E0-3
friction_rule: >
  Drake combines a pair as 2 mu_a mu_b / (mu_a + mu_b).  The object geometry carries OBJECT_MU and
  scene.partner_mu gives each strip and the paddle the value that makes the PAIR coefficient equal to
  its target, so one object can touch strips of different friction.
\end{verbatim}

\end{document}
